\section*{Chapitre \ref{ch:b3:x-ray-diffraction} --- Cristallographie et diffraction des rayons X}

\begin{solution}{exo:b3:x-ray-diffraction:1}
$d = a/\sqrt N$~: $d_{111} = \qty{325.66}{pm}$, $d_{200} = \qty{282.03}{pm}$, $d_{220} =
\qty{199.42}{pm}$.
\end{solution}

\begin{solution}{exo:b3:x-ray-diffraction:2}
Cubique à faces centrées~: les premières réflexions sont (111) et (200).
$\sin\theta = \lambda\sqrt N/2a$~: $2\theta =
\ang{43.32}$ et \ang{50.45}.
\end{solution}

\begin{solution}{exo:b3:x-ray-diffraction:3}
P~: toutes les sept. I ($h + k + l$ pair)~: 110, 200, 211, 220. F (indices non mélangés)~: 111, 200, 220.
\end{solution}

\begin{solution}{exo:b3:x-ray-diffraction:4}
$\mathbf a^* = \mathbf e_x/a$, $\mathbf b^* = \mathbf e_y/b$~: un réseau
rectangulaire de côtés $1/a$
et $1/b$, le grand côté du réseau direct devenant le petit.
\end{solution}

\begin{solution}{exo:b3:x-ray-diffraction:5}
Rapports des $\sin^2\theta$~: $1 : 2 : 3 : 4 : 5$, c'est-à-dire
$N = 2, 4, 6, 8, 10$~: réseau centré (110,
200, 211, 220, 310). D'après la première raie, $d_{110} = \lambda/2\sin(\ang{20.135}) =
\qty{223.77}{pm}$ et $a = d\sqrt2 = \qty{316.5}{pm}$~: le tungstène.
\end{solution}

\begin{solution}{exo:b3:x-ray-diffraction:6}
$Z = \rho N_Aa^3/M = 1.99 \times 6.022\times10^{23} \times (6.2929\times10^{-8})^3/74.6 = 4.0$.
\end{solution}

\begin{solution}{exo:b3:x-ray-diffraction:7}
$\beta = \qty{0.40}{\degree} = \qty{6.98e-3}{rad}$, $\cos\theta = \cos\ang{15.85} = 0.962$~:
$L = 0.9 \times 0.15406/(6.98\times10^{-3} \times 0.962) = \qty{21}{nm}$.
\end{solution}

\begin{solution}{exo:b3:x-ray-diffraction:8}
$F_{hkl} = f_{\ce{Cs+}} + f_{\ce{Cl-}}(-1)^{h+k+l}$~: $F_{100} = 54 - 18 = 36$, $F_{110} = 72$~:
intensités dans le rapport 1~:~4. Le centre est occupé par un ion
différent de celui des sommets~: le réseau est primitif (un réseau centré
exige le même contenu aux deux points), et la réflexion 100 est
présente.
\end{solution}

\begin{solution}{exo:b3:x-ray-diffraction:9}
Chaque atome a une copie décalée de $(\frac12,\frac12,0)$, ce qui
multiplie $F$ par $1 +
\eu^{\iu\pi(h+k)} = 1 + (-1)^{h+k}$, nul pour $h + k$ impair.
\end{solution}

\begin{solution}{exo:b3:x-ray-diffraction:10}
Les quasicristaux présentent un ordre à longue distance sans
périodicité~: ils n'ont pas de réseau, si bien que la restriction, qui
porte sur les réseaux, ne s'applique pas à eux. Leurs taches de
diffraction nettes ont conduit à définir un cristal comme tout solide
dont le diagramme de diffraction est essentiellement discret, périodique
ou non.
\end{solution}

\begin{solution}{exo:b3:x-ray-diffraction:11}
Le glissement $c$ envoie $(x, y, z)$ en $(x, \frac12 - y, \frac12 + z)$.
Pour $k = 0$, les deux
facteurs de phase sont $\eu^{2\pi\iu(hx + lz)}$ et $\eu^{2\pi\iu(hx + lz + l/2)} = (-1)^l
\eu^{2\pi\iu(hx+lz)}$~: leur somme s'annule pour $l$ impair.
\end{solution}

\begin{solution}{exo:b3:x-ray-diffraction:12}
Une inversion, un miroir ou un glissement transforme une molécule en son
image dans un miroir, si bien qu'un cristal qui en contiendrait un
contiendrait les deux énantiomères. Un énantiomère pur ne cristallise que
dans les 65 \omterm{def:b3:x-ray-diffraction:space-group}{groupes d'espace} construits à partir de rotations, d'axes
hélicoïdaux et de translations (groupes de Sohncke), comme
\textit{P}$2_12_12_1$. D'après la loi de Friedel, le diagramme d'un
énantiomère est égal à celui de l'autre~; la configuration absolue
s'obtient à partir des petits écarts causés par la diffusion anomale
d'atomes plus lourds.
\end{solution}

\begin{solution}{pb:b3:x-ray-diffraction:1}
\textbf{1.} $\theta = \ang{14.171}$, $d = \lambda/2\sin\theta = \qty{314.64}{pm}$.
\textbf{2.} 314.64, 222.49, 181.66, 157.32, 140.71, \qty{128.45}{pm}.
\textbf{3.} 1.000, 2.000, 3.000, 4.000, 5.000, 6.000.
\textbf{4.} 1, 2, 3, 4, 5, 6.
\textbf{5.} Oui, comme les réflexions (100), (110), (111), (200), (210),
(211) d'une maille cubique primitive
d'arête $a' = \qty{314.6}{pm}$.
\textbf{6.} $N = 7$ (ainsi que 15, 23\dots), qui n'est pas une somme de
trois carrés~; la première différence apparaîtrait à la huitième raie.
Six raies ne permettent pas de distinguer une maille P d'arête $a'$
d'une maille F d'arête $2a'$.
\textbf{7.} $N = 4, 8, 12, 16, 20, 24$~: (200), (220), (222), (400), (420), (422).
\textbf{8.} $a = 2d_{200} = \qty{629.3}{pm}$.
\textbf{9.} \ce{KCl} (\qty{629.29}{pm}).
\textbf{10.} $F = 4[f(\ce{K+}) - f(\ce{Cl-})]$ pour des indices tous
impairs, et les deux ions ont chacun 18
électrons~: les amplitudes se compensent.
\textbf{11.} Oui pour les deux~: \ce{Na+} (10) et \ce{Cl-} (18), \ce{K+}
(18) et \ce{Br-} (36)
diffusent différemment~; la raie (111) de \ce{NaCl} serait à
\ang{27.36}, celle de \ce{KBr} à \ang{23.38},
toutes deux dans le domaine enregistré.
\textbf{12.} Les rayons X voient \ce{K+} et \ce{Cl-} comme presque
identiques~; des ions identiques disposés en sel gemme forment un
arrangement cubique simple d'arête $a/2$, une maille apparente.
\textbf{13.} Quatre \ce{KCl}.
\textbf{14.} $\rho = 4 \times 74.6/(6.022\times10^{23} \times (6.293\times10^{-8})^3) =
\qty{1.99}{g/cm^3}$.
\textbf{15.} Six et six (octaèdres de l'autre ion).
\textbf{16.} $a/2 = \qty{314.6}{pm}$.
\textbf{17.} (440), $N = 32$, à \ang{87.65} (les réflexions à indices
tous impairs (333)/(511) étant absentes).
\textbf{18.} Les intensités varient avec l'orientation préférentielle des
cristallites, l'absorption, la décroissance des facteurs de diffusion et
l'agitation thermique~; les positions ne dépendent que du réseau.
\textbf{19.} D'après $d = \lambda/2\sin\theta$, $\Delta d/d = -\cot\theta\,\Delta\theta$~:
une même erreur angulaire donne une erreur relative plus petite aux
grands $\theta$.
\textbf{20.} $d_{420} = \lambda/2\sin\ang{33.190} = \qty{140.71}{pm}$, $a = d\sqrt{20} =
\qty{629.3}{pm}$.
\textbf{21.} $L = 0.9 \times 0.15406/(4.36\times10^{-3} \times \cos\ang{33.19}) = \qty{38}{nm}$.
\textbf{22.} Non~: $\beta = \ang{0.0095}$, bien au-dessous d'une largeur
instrumentale typique.
\textbf{23.} \textbf{$a = \qty{629.3}{pm}$~: la poudre est du chlorure de
potassium}, dont les réflexions impaires sont éteintes par deux ions
ayant le même nombre d'électrons.
\end{solution}
